\paragraph{The tests on the three level rings.}
\label{par:fitted-arc}
\hypertarget{source-ordinary-marked-points}{}
The curves lie in a frame. Their specializations at the three levels must pass the
tests of Section~\ref{sec:marked-patching}. In this
paragraph \(P=\bar\Psi\) denotes the character of the fitted field
construction in Section~\ref{sec:marked-hosts}. The residual characters
on the CM group at the level are
\(\bar\kappa_n P^{w_i}\), with \(w=(1,-1)\) or \((1,0,-1)\).
The point on the level is the pullback of the chart representation
by the surjection to the frame group, multiplied by \(\kappa_n\).
Its ordinary characters are multiplied by the same scalar. Thus the
specialized determinant and ordered characters evaluate by the map
from the actual level ring to the chart followed by the map to the
curve DVR. In particular its kernel is the required contracted prime.
The finite scalar and the cyclotomic twist have the prescribed
polarization, as in Section~\ref{sec:linked-seeds}.

\begin{auxlemma}[Multiplicity of residue characters]\label{lem:residue-character-multiplicities}
Let \(k\) be a field of characteristic \(\ell>0\), and let \(\chi\) range over a finite set of distinct characters \(H\to k^\times\). If two lists of nonnegative integer multiplicities satisfy
\[
 \prod_\chi (T-\chi(h))^{a_\chi}
       =\prod_\chi (T-\chi(h))^{b_\chi}\quad\text{in }k[T]
       \qquad(h\in H),
\]
then \(a_\chi=b_\chi\) for every \(\chi\).
\end{auxlemma}
\begin{proof}
The coefficient below the leading term gives \(\sum_\chi (a_\chi-b_\chi)\chi=0\) as functions on \(H\), with the differences read in \(k\). Distinct characters are linearly independent, so \(a_\chi-b_\chi\in\ell\mathbb Z\). Cancel the \(\min(a_\chi,b_\chi)\) common factors for each \(\chi\). The remaining multiplicities on both sides are multiples of \(\ell\), giving
\[
 \left(\prod_\chi(T-\chi(h))^{(a_\chi-b_\chi)_+/\ell}\right)^\ell
 =
 \left(\prod_\chi(T-\chi(h))^{(b_\chi-a_\chi)_+/\ell}\right)^\ell ,
 \qquad u_+=\max(u,0).
\]
Frobenius is injective on \(k[T]\), so the two inner polynomials agree at each \(h\). The remaining total degree strictly decreases at every iteration that starts in positive degree. It reaches zero, proving equality of all original multiplicities.
\end{proof}

\begin{proposition}[Local labels and spans at a fitted point]
\label{prop:fitted-points-three-tests}
At the indicated points of the three level rings the ordered labels
remain distinct on the specified local groups after restriction
to an open subgroup. Write \(H\) for the full cyclotomic kernel on the positive component.
Over an algebraic closure of the coefficient field the two block spans are
\[
 \operatorname{span}r(H)=
 \begin{cases}
   \operatorname{End}(V),&\text{type I},\\
   \operatorname{End}(V_{\rm out})\oplus\operatorname{End}(V_{\rm mid}),&\text{type E},
 \end{cases}
\]
where \(V_{\rm out}\) is the outer plane and \(V_{\rm mid}\) the middle line. In the
open-coset test the element with separate eigenvalues is chosen
in that coset.
\end{proposition}

\begin{proof}
\proofstep{Testing the ordered local labels.}
At each \(v\mid\ell\), on the preimage of the distinguished local
subgroup \(N_v\) the residual character ratios are one and the full
cyclotomic character is one. The former follows from \(P|_{N_v}=1\),
and the latter is part of the definition of \(N_v\).
The nice curve avoids every local coincidence ideal, so each
distinct pair of labels has a nonidentity ratio somewhere on this
subgroup. This persists under the scalar twist. The group
\(1+t k[[t]]\) is torsion-free in characteristic \(\ell\):
raising to \(\ell^a\) raises the degree of the leading nonconstant
term by that factor, and raising to a power prime to \(\ell\)
multiplies its coefficient by a unit.

For a ratio \(\alpha\) nontrivial on an element \(n\) in
the preimage of \(N_v\), choose \(m>0\) with \(n^m\) in any
prescribed open subgroup of the global group. Then
\[
    \alpha(n^m)=\alpha(n)^m\ne1,\qquad
    \epsilon(n^m)=1 .
\]
Thus the same pair of labels is distinct on the restricted
cyclotomic kernel. We can choose a single element there on which
all ratios have infinite order. To see this, successively multiply by
powers of elements on which the individual ratios are nontrivial. A previously nontrivial
ratio can become one for at most one exponent if it is nontrivial
on the new element, and otherwise stays unchanged. The new ratio
can also be made nontrivial by avoiding one exponent. This supplies
an element of the local group in the chosen open subgroup and the
full cyclotomic kernel with pairwise distinct eigenvalues at
\emph{every} positive power. At each marked place the ordering must pass both tests simultaneously:
\[
 \psi_i\ne\psi_j,\qquad \psi_i/\psi_j\ne\epsilon\qquad(i<j).
\]

\proofstep{Full spans on the cyclotomic kernel.}
We next check absolute irreducibility on small subgroups of the
cyclotomic kernel. Write \(G\) for the level CM group and
\(H=\ker(\epsilon|_G)\). Work over an algebraic closure of the
fraction field of the curve. In type I the marked representation is
irreducible on the frame group, hence on \(G\) after pullback and
scalar twist. In type E apply this assertion to its outer plane,
irreducible since neither outer triangular ideal vanishes. Its two
residual character labels are the outer labels above. Indeed on the
residual triangular chart the diagonal entries give those characters
and the product of the two opposite outer entries is zero.

On the outer plane the ordinary labels are obtained by removing the middle character with multiplicity one. If \(S\) denotes the multiset for the full flag, then
\[
 S=S_{\rm outer}\sqcup\{\chi_{\rm middle}\}.
\] One way to see the deletion
is to take a local element with three distinct eigenvalues. Both
summands and all flag steps are sums of its eigenlines, so their
intersections give the induced flags.

In either irreducible block the dimension \(b\) is a prime, two or
three. Let \(U\) be any open subgroup of \(G\) and take
\(H'=H\cap\operatorname{core}_G(U)\), normal in \(G\).
If \(W_0\) is a simple \(H'\)-submodule, every \(gW_0\)
is again simple for \(H'\) by normality. Their sum is a
nonzero \(G\)-submodule, hence the full irreducible block.
This proves semisimplicity. Irreducibility for \(G\) makes
the action on \(H'\)-isotypic pieces transitive. If \(t\)
counts the pieces, \(d_1\) is the dimension of a simple
piece's constituent and \(m_1\) its multiplicity, then
\[
                             b=t\,d_1\,m_1 .
\]
Since \(b\) is prime, reducibility forces either \(t=d_1=1\)
(scalar action) or \(d_1=m_1=1\) and \(t=b\) (distinct
character lines permuted transitively).
The separating local element in \(H'\) excludes scalars.

In the
second case let \(G_1\) be the kernel of the permutation action.
On the intersection of the local decomposition group with \(H'\),
the characters of the lines agree up to permutation with the
restrictions of the ordinary labels. This follows by comparing
the semisimplification of the flag with the direct sum of the lines.
Write \(D\) for the local decomposition group, and \(\lambda_i\) for the distinct characters on \(D\cap H'\). Each is the restriction of a character \(\widetilde\lambda_i\) of \(D\); hence for \(d\in D\) and \(h\in D\cap H'\),
\[
 \lambda_i(dhd^{-1})=\widetilde\lambda_i(d)\lambda_i(h)\widetilde\lambda_i(d)^{-1}
 =\lambda_i(h).
\]
The corresponding eigenspaces are distinct lines. Every \(d\) preserves each line, so \(D\subseteq G_1\).

On \(G_1\) the line characters are units of the valuation ring of the
algebraic closure: their values and inverses are roots of monic
characteristic polynomials over the curve DVR. Reducing gives
characters whose product characteristic polynomial at each element
is that of the prescribed global residual labels restricted to
\(G_1\). The values of the reduced line characters belong to the finite image of the residual diagonal characters. The two lists consequently factor through a common finite abelian quotient of order prime to \(\ell\). Applying Lemma~\ref{lem:residue-character-multiplicities} on \(G_1\) identifies their multisets with their full integer multiplicities. Write \(\chi\) for a reduced character on \(G_1\) and \(\tilde\chi\) for a global extension. Then conjugation fixes \(\chi\), since
\[
 \chi(ghg^{-1})=\tilde\chi(g)\chi(h)\tilde\chi(g)^{-1}=\chi(h)\qquad(h\in G_1).
\] Transitivity of the original lines forces their
reductions to agree on \(G_1\), including the local group.

The local
mark with \(P(g_*)^2\ne1\) distinguishes the outer residual labels
there. This contradiction proves irreducibility on \(H'\), hence on
\(U\cap H\). The local mark exists at each place by conjugating
the mark of \eqref{eq:tag-h1} in the abelian fitted field and using the
global character \(P\). Equivalently one place suffices for this
contradiction. In type E the two simple \(H\)-constituents are inequivalent
because their dimensions are different. The semisimple image
algebra of the group ring acts as the product of their full
endomorphism algebras by density. In the two cases the spans
over the algebraic coefficient field are therefore
\[
  \operatorname{span}r(H)=
   \begin{cases}
      \operatorname{End}(V),&\text{type I},\\
      \operatorname{End}(V_{\rm out})\oplus\operatorname{End}(V_{\rm mid}),
          &\text{type E}.
   \end{cases}
\]
The equality follows over the original coefficient field by
extension of scalars, since the absolute dimensions are the
displayed block dimensions.

\paragraph{Regular elements in each open coset.} Both block patterns now require the same density statement, with the group element
constrained to a given open coset. In an open coset meeting \(H\)
choose \(s\in H\) and an open subgroup \(U\) with
\(s(U\cap H)\) inside the coset. Let \(C\) be the Zariski closure
of the image of \(U\cap H\) and \(C^\circ\) its identity component
over the algebraic closure. The preimage of \(C^\circ\) is open in
\(U\cap H\). Indeed each of finitely many components is closed for
the adic topology on matrix entries: individual polynomial equations
over the algebraic closure already have coefficients in a finite
extension. Choose an open subgroup \(V\subset U\) of \(G\) small enough that 
\[
 V\cap H\subseteq r^{-1}(C^\circ),\qquad r(V\cap H)\subseteq C^\circ.
\]
 The irreducible blocks just proved
are irreducible for \(C^\circ\).

\proofstep{A diagonal torus in the connected image.}
Take the separating local element in \(V\cap H\) and diagonalize its
matrix. The identity component of the closure of its powers is a
torus inside the diagonal torus of \(C^\circ\). Let \(x\) be the matrix of this element and let \(T\) be the identity component of the closure of its powers. For distinct eigencharacters \(\lambda_i,\lambda_j\), coincidence on \(T\) would yield 
\[
 (\lambda_i/\lambda_j)(x)^m=1,\qquad m=[\overline{\langle x\rangle}^{\rm Zar}:T],
\]
 contrary to the separating property. In the rational cocharacter space the equal-coordinate loci are finitely many proper hyperplanes. Choose a rational point in their complement and clear denominators to obtain the integral cocharacter
\[
 d=(d_i),\qquad \prod_{i<j}(d_i-d_j)\ne0 .
\]

In this eigenbasis there is \(h\in C^\circ\) such that all principal
minors of \(A=r(s)h\) are nonzero. For one block of dimension at
most three the first and the next-to-last exterior powers are
irreducible: the latter is the dual of the first tensored by the
determinant. The other nonempty exterior power is the determinant
itself. For each principal minor, the orbit of its basis wedge
therefore spans the corresponding exterior power.

For a wedge \(e_I\) and its coordinate functional \(e_I^*\), the principal minor equals 
\[
 \det (r(s)h)_{I,I}=e_I^*\!\left(\bigl(\textstyle\bigwedge^{|I|}(r(s)h)\bigr)e_I\right).
\]
 The orbit spans and the functional composed with \(\bigwedge^{|I|}r(s)\) is nonzero, so this minor is nonzero somewhere on \(C^\circ\).
For the outer plane plus line, a principal minor is the product of
a principal minor on the plane and, when the line is included, its
nonzero scalar. The same argument on the plane works. The simultaneous
choice follows from irreducibility of the connected algebraic group.

\proofstep{A separable characteristic polynomial in the desired coset.}
Consider \(A\operatorname{diag}(z^{d_i})\) over the Laurent
series field in \(z\). Its characteristic coefficient at \(X^{n-j}\)
is
\[
 (-1)^j\sum_{|I|=j}\det A_{I,I}\,z^{\sum_{i\in I} d_i}.
\]
The subset of the \(j\) smallest exponents uniquely minimizes
the order and its principal minor is nonzero. Number the ordered
exponents \(d_{(1)}<\cdots<d_{(n)}\). The consecutive coefficient
orders, read from the leading term, differ by the successive
\(d_{(i)}\). The Newton polygon therefore has \(n\) segments of
horizontal length one with distinct slopes. Its roots
over the algebraic closure have distinct valuations. Hence the
discriminant on \(r(s)C^\circ\) is a nonzero regular function, since
its pullback along this cocharacter translate is nonzero.
The image of \(U\cap H\) is Zariski dense in \(C\), so some actual
element of \(s(U\cap H)\) has nonzero discriminant. This proves the
required density in the profinite topology, including cosets outside
the identity component.

\end{proof}

\paragraph{Integral normalization and the distinguished place.}
To apply the tests at the distinguished place we need integral representatives of the
blocks. After a finite coefficient extension use coordinates on the blocks
defined over the fraction field of the new DVR. In type E their
dimensions differ, so polarization pairs each block with itself:
a cross pairing would intertwine inequivalent simple modules after
conjugation and duality. The form on each block is nondegenerate
and symmetric.

In a block \(V_i\) over a coefficient extension \(E'/E\), intersect the extended lattice with that subspace: 
\[
 L_i=(L\otimes_{O_E}O_{E'})\cap V_i,\qquad r(g)L_i=L_i.
\]
 Here \(L\) is the original stable lattice and \(O_{E'}\) the extended coefficient DVR. The dual is stable as well. Indeed, the polarized action transports integral pairings to integral pairings, since its multiplier is a unit and conjugation preserves the group. Applying the same argument to the inverse gives, for the dual lattice \(\Lambda^\vee\),
\[
 r(g)\Lambda^\vee\subseteq\Lambda^\vee,\quad r(g^{-1})\Lambda^\vee\subseteq\Lambda^\vee\quad\Longrightarrow\quad r(g)\Lambda^\vee=\Lambda^\vee .
\] Diagonalize the Gram matrix over the DVR,
allowing fractional entries. One may first scale to an integral
matrix and choose a primitive vector with self-pairing of smallest
valuation among the entries. When a mixed entry has that valuation
and both diagonal entries have greater valuation, take the sum of
the two basis vectors, using \(2\) invertible. Split off this
vector by integral row and column operations and continue.

Write the diagonal entries as \(t^{a_i}u_i\) with \(u_i\) units.
In this basis preservation of the lattice and its dual requires
\[
       v_t(r_{ij})\ \ge\ \max(0,a_j-a_i).
\]
After adjoining \(\sqrt t\), use the basis \(e'_i=t^{-a_i/2}e_i\).
In the extended valuation with \(v_t(t)=1\), the new entries satisfy
\[
  v_t(r'_{ij})=v_t(r_{ij})+\tfrac12(a_i-a_j)
                 \ \ge\ \tfrac12|a_i-a_j|.
\]
The same calculation applies to inverse matrices, so the new lattice
is stable. Its Gram matrix is \(\operatorname{diag}(u_i)\).
Adjoin square roots of the reductions of the \(u_i\) by a finite
residue extension and lift them by Hensel's lemma. Put \(e_i=u_i^{-1/2}e'_i\). Since the chosen square roots are units, this preserves the lattice and normalizes each diagonal Gram entry:
\[
 (u_i^{-1/2})^2u_i=1 .
\]

In each ordinary flag a local matrix with distinct
eigenvalues splits its steps into eigenlines over the fraction
field. These eigenlines lie in the separate blocks.
Intersect the resulting flags in each block with its lattice.
The steps are saturated, and successive free quotients over the
DVR give an integral basis compatible with their order. Interleave
the block bases in the prescribed local order. The original and new
coordinates are related over a finite extension by fixed matrices,
so adic continuity, the two character inequalities, and the span
and separability tests persist. The coefficient maps are injective
and local, hence keep the evaluation kernels and determinant
identities. This supplies the normalized points used for the
simultaneous prime choice.

\begin{lemma}[The marked support tests on the chosen curves]
\label{lem:fitted-marked-support-test}
In the integral normalization the distinguished-place test agrees
with the auxiliary level maps of the three rings. Let \(b\) be the trace Gram minor nonzero at the type-I point \(Q\).
The same minor detects the full matrix span at its generizations since
\[
           P\subseteq Q,\quad b\notin Q\quad\Longrightarrow\quad b\notin P.
\]
The support test therefore applies at the indicated type-I generizations;
the two-block and augmentation comparisons give support at the
indicated stable or thick type-E target from the stable start.
\end{lemma}

\begin{proof}
\proofstep{The distinguished-place vanishing test.}
There is a split distinguished place for which the twisted dual
test holds for \emph{all} these curves. On the field \eqref{eq:tag-h2} choose
a CM element with \(P=1\) and \(\bar\epsilon\ne1\). Explicitly
multiply complex conjugation by the element with \(P=\bar\epsilon=1\),
\(\eta=-1\) in Section~\ref{sec:marked-hosts}. If necessary multiply
by inertia at the new prime ramified in \(\Q(\sqrt{Dp})\) to enter
\(H_0\); the three characters \(P,\bar\epsilon,\eta\) are unramified
there. Raise to a common \(\ell\)-power to enter the three
\(\ell\)-power kernels of \eqref{eq:tag-h2}.

If \(g\) is the adjusted element, the common \(\ell\)-power preserves its cyclotomic residue because \(\ell\) is odd: 
\[
 \bar\epsilon(g^{\ell^a})=\bar\epsilon(g)^{\ell^a}=(-1)^{\ell^a}=-1.
\]
 Chebotarev over \(L^+\) for the finite quotient recording
\(P,\bar\epsilon,\eta\) gives places \(v_0\) split in \(L\) with
such Frobenius, away from the cutoff and the discriminant. One can require residue degree one over \(\mathbb Q\). Indeed the higher-degree primes of a fixed number field \(F\) contribute a bounded sum as \(s\downarrow1\), since
\[
 \sum_{\mathfrak p:\ f(\mathfrak p/p)>1}(N\mathfrak p)^{-s}\le [F:\mathbb Q]\sum_p p^{-2s}=O_F(1).
\] Choose the second omitted place away from
the first. At each chosen decomposition group all the twisted
residual characters have the common value \(\bar\kappa_n\).

For a Frobenius matrix \(T\) at an arc and cyclotomic value \(e_0\),
the reductions of \(T\) and \(e_0 T\) have single, distinct,
nonzero eigenvalues. Their characteristic polynomials have unit
resultant over the curve DVR. Let \(P,Q\) be the characteristic polynomials of
\(e_0 T,T\) over the curve DVR. Since the resultant is a unit, the two polynomials generate the unit ideal, giving a B\'ezout identity over the coefficient ring:
\[
 AP+BQ=1 .
\]
If \(e_0TY=YT\), then \(f(e_0T)Y=Yf(T)\) for every
polynomial \(f\). Applying the Bézout relation on the
left gives
\[
  Y=A(e_0T)P(e_0T)Y+B(e_0T)YQ(T)=0
\]
by Cayley--Hamilton. This only needs the reduced
characteristic polynomials, so it covers nonsemisimple
reduction. This gives the distinguished-place vanishing test and the nonidentity cyclotomic scalar used in the paired detector.

At each of the three levels take the finite set consisting of all
places above \(\ell\), the tame places above the cutoff away from
\(\ell\), and its own distinguished \(v_0\). All split in the CM
extension, which is unramified at finite places. The base Galois
level has exactly this set of allowed ramification. At the
distinguished place use the first principal congruence subgroup,
at the fixed tame places the coefficient type for the residual
inertia characters, and at \(\ell\) the ordered ordinary torus
character, with the twists of Section~\ref{sec:linked-seeds}.

If \(u\) has order dividing \(\ell^a\), its value under a local character \(\theta\) into the curve DVR satisfies
\[
 (\theta(u)-1)^{\ell^a}=\theta(u)^{\ell^a}-1=0.
\]
The DVR is reduced, so \(\theta(u)=1\). Under local reciprocity at a split place \(w\mid\ell\), the free pro-\(\ell\) unit rank is \([L_w:\Q_\ell]\); a uniformizer supplies the independent unramified generator.
The base evaluation kernel on the empty-auxiliary Hecke action is
precisely \(J_2,J_3\) or \(J_0\) according to the chosen ring.

\proofstep{Compatibility with the auxiliary level maps.}
For each selected system of \(q\) further places use the two maps
of Section~\ref{par:ordinary-level} for that same ordinary module.
Their composites are the same split-root factor, the tame
augmentation is given by the trace, and the relation bound there is
\(n!\) at every auxiliary depth. The proofs use rank \(n\) compact
and Iwahori fixed vectors, regular ordinary specialization and the
fixed tame type. At the two marked levels the relation bound \(n!\) reads
\[
 2!=2\quad(n=2),\qquad 3!=6\quad(n=3)
\] Choose one split completion \(w\) above each \(v\mid\ell\) of \(L^+\), so \(L_w\simeq L^+_v\). The \(n\) ordinary unit characters therefore give
\[
 n\sum_{v\mid\ell}[L^+_v:\Q_\ell]=n[L^+:\Q]
 \quad(n=2\text{ or }3).
\]
The two rank-three actions are compared to their own base ideals before transporting either to the common frame.

The ordinary deformation and the ordered characters of each curve
are defined by the same pullback and scalar as the map \(\Phi_n\).
Let \(Q\) be the type-I curve kernel and \(b\) its nonzero trace Gram minor. Every generization \(P\) satisfies 
\[
 P\subseteq Q,\quad b\notin Q\quad\Longrightarrow\quad b\notin P.
\]
 Starting at that point in the Hecke support,
the one-block local comparison and augmentation put every such
generization in the support.

At a chain of type-E points the two
blocks above give the hypotheses of the two-block local comparison.
Sections~\ref{subsec:closed-link} and~\ref{subsec:marked-augmentation}
then give support on the stable or thick target from a stable
start in the support. These are the prime conditions at the three
marked ideals under the ordinary level construction of this
supplement. Constants for curve normalization, cohomology and the
number of tests may depend on the fitted field, level and finite
list of curves. The patching conclusion is ideal containment and
introduces no such factor in the length estimate on the original field. \label{cor:ordinary-marked-points-1}
\end{proof}
